\documentclass{upmassignment}
\usepackage[spanish]{babel}
\usepackage{ifthen}
\usepackage{amsmath}
\usepackage{amsfonts}



% Para mostrar/ocultar soluciones
\newboolean{show}
\setboolean{show}{true}
%\setboolean{show}{false}
\usepackage{environ}
\NewEnviron{solucion}{
  \ifshow
      \begin{answer}\BODY\end{answer}
  \fi}






\coursetitle{Creating assignments}
\courselabel{RSTC}
\exercisesheet{Ejercicio de Alumno}{Tema 5}
\student{\ }%
\semester{Segundo Semestre}
\date{\today}
\university{Universidad Politécnica de Madrid}
\school{Departamento de Ingeniería de Sistemas Telemáticos}
%\usepackage[pdftex]{graphicx}
%\usepackage{subfigure}


\setlength{\textwidth}{5.0in}
\linespread{1.3}
\renewcommand{\PB}{{\bfseries Problema}}















\begin{document}


Considere un switch con tamaño de cola infinita
al que llega tráfico Poissoniano de tasa
$\lambda$. Sabiendo que el tiempo de
servicio del switch sigue una distribución
uniforme $U(a,\tfrac{2}{\mu}-a)$, responda
justificadamente a las siguientes preguntas:

\begin{problemlist}
    \pbitem Obtenga el tiempo medio
        de tránsito por el sistema.

    \begin{solucion}
        \input{../../RSTC-solutions/tema-6/solucion2026-problema1.tex}
    \end{solucion}


    \pbitem Encuentre el valor de $a$
        que minimiza el tiempo medio residual.
        \begin{solucion}
            \input{../../RSTC-solutions/tema-6/solucion2026-problema2.tex}
        \end{solucion}

    \pbitem Suponga que el tiempo de servicio
        del switch es exponencial de tasa
        $\mu$. El $p$ porciento del
        tráfico saliente va a otro switch
        idéntico con tasa de servicio $2\mu$,
        i el $1-p$ porciento restante sale
        del sistema.
        ¿Cuál es el tiempo medio de
        tránsito de un paquete?
        \begin{solucion}
            \input{../../RSTC-solutions/tema-6/solucion2026-problema3.tex}
        \end{solucion}

\end{problemlist}

\end{document}


